\section{The dimension estimate and the choice of parameters}
\label{sec:conclusion}

We can now assemble the argument. Section~\ref{sec:arithmetic} made the
forms integral, Section~\ref{sec:estimates} bounded their coefficients
and values, and Section~\ref{sec:rank} supplied the rank statement.
We first apply the criterion with all auxiliary parameters fixed, then
choose them as the number of beta values increases.

For each admissible $n$, form the integer matrix with rows the coefficients
of $\Lambda_{n,p,k}$ in
\[
 1,\beta(2),\ldots,\beta(2\lfloor(a+h)/2\rfloor).
\]
It satisfies \eqref{eq:weak-rank}.
Indeed, a rational vector annihilating all these rows lifts to one for
the rows $Q_{k,i}^{[p]}(-1)$ by putting zero in every odd coordinate and
multiplying coordinate $i/2$ by $-2^{i+1}$ for even $i\geq2$.
Its constant coordinate is unchanged.
The row multipliers $\Dt_{d,k-1}/(k-1)!$ are nonzero, so
Proposition~\ref{prop:rank} forces that constant coordinate to be zero.

If $\log\alpha<0$, Proposition~\ref{prop:criterion} gives
\[
 \dim_\Q\Span_\Q\{1,\beta(2),\ldots,
                      \beta(2\lfloor(a+h)/2\rfloor)\}
       \geq1-\frac{\log\alpha}{\log B}.
\]
Removing $1$ loses at most one dimension. Thus
\begin{equation}\label{eq:parameter-bound}
 D(2\lfloor(a+h)/2\rfloor)\geq-\frac{\log\alpha}{\log B}.
\end{equation}

The scales in the parameter choice can be read from this quotient.
Increasing $\Omega$ improves the decay through $\Omega\log r$, but it
also increases $\log\chi$ by a term of order $\Omega^2/a$.
The denominator estimate contributes a term of order $\log a$ when
$h$ is proportional to $a$. Taking $\Omega$ of order
$\sqrt{a\log a}$ therefore keeps $\log B$ of order $\log a$ while
making $-\log\alpha$ of order $\sqrt{a\log a}$. Their quotient has the
desired order $\sqrt{a/\log a}$.

For the numerical constants, use the choice of \cite[Section~4.6]{F26}.
Let $a$ tend to infinity through positive multiples of $25$, and set
\begin{equation}\label{eq:parameters}
 h=\frac{9a}{25},\quad
 r=\frac{39}{10},\quad
 \kappa=\frac{529}{50},\quad
 \omega=12,\quad
 \Omega=\left\lfloor r\sqrt{a\log a}\right\rfloor.
\end{equation}
For sufficiently large $a$ these satisfy
$a>\Omega\geq\omega>\kappa>2r$ and $0\leq h\leq a$, and
\[
 (h+1)(\kappa-2r)+\omega
  =\left(\frac{9a}{25}+1\right)\frac{139}{50}+12>a.
\]
Take $n$ through sufficiently large multiples of $50$.
For each fixed $a$, all needed integer conditions and the weak rank
statement then hold.

As $a\to\infty$, equations \eqref{eq:coeff-small}, \eqref{eq:B} and
\eqref{eq:alpha} give
\begin{align}
 \log\chi&\sim\frac{r^2\log r}{2}\log a,\notag\\
 \log B&\sim\left(\frac{r^2\log r}{2}+\kappa\right)\log a,
                                                   \label{eq:rates-asymptotic}\\
 \log\alpha&\sim-r\log r\sqrt{a\log a}.\notag
\end{align}
In particular, $\log\alpha<0$ for sufficiently large $a$.
With $s=a+h=34a/25$, \eqref{eq:parameter-bound} implies
\begin{equation}\label{eq:cstar}
 \liminf_{\substack{s\to\infty\\34\mid s}}
   \frac{D(s)}{\sqrt{s/\log s}}
 \geq
 c_*:=\frac{r\log r}
   {\left(r^2\log r/2+\kappa\right)\sqrt{34/25}}
 =0.2174562348\ldots .
\end{equation}
The strict inequality $c_*>0.21$ requires little numerical precision:
$1.36<\log(3.9)$ and $\sqrt{1.36}<1.167$ already suffice,
since $x/(r^2x/2+\kappa)$ increases with $x>0$.

Finally $D(s)$ is nondecreasing, and successive multiples of $34$ have
bounded gaps. For an arbitrary large even $s$, apply
\eqref{eq:cstar} to $34\lfloor s/34\rfloor$.
The quotient of the corresponding factors $\sqrt{s/\log s}$ tends to $1$.
The strict slack between $c_*$ and $0.21$ proves
Theorem~\ref{thm:main}.

\begin{remark}
The limiting assertion \eqref{eq:cstar} extends to all even $s$ in the
same way. We retain the simpler constant $0.21$ in the theorem.
No explicit threshold for \emph{sufficiently large} is asserted;
in particular, the multiplicity constant has not been tracked.
\end{remark}

\section*{Acknowledgements and preparation}

The proof uses Fischler's rational-function construction, refined
linear-independence criterion and multiplicity method. The numerical
parameters and the organization of the rank argument are taken from his
proof, with the adaptations identified in the text.
Mistiaen's explanation of the conductor restriction suggested the
evaluation-point question addressed here.

This research note was prepared with OpenAI Codex. The accompanying exact
calculations check finite instances of the identities and denominator
bounds; they are not formal verification. The paper has been independently
cross-checked by another model.

